\RequirePackage{fix-cm}
\documentclass[12pt,letterpaper]{article}
\usepackage[margin=.5in]{geometry}
\usepackage{tikz}
\usetikzlibrary{arrows.meta,patterns}
\pdfmapfile{+cm.map}
\pagestyle{empty}
\setlength{\parindent}{0pt}
\begin{document}
\null\begin{tikzpicture}[remember picture,overlay]
\begin{scope}[shift={(current page.north west)},x=1cm,y=-1cm]
\begin{scope}[shift={(1.5,1.5)}]
\node[anchor=north west,inner sep=0,align=left,text width=18.5cm,font=\sffamily\fontsize{11}{13.75}\selectfont] at (0,-0.5) {Week 35 / Footprint borders / K--1};
\node[anchor=north west,inner sep=0,align=left,text width=17cm,font=\sffamily\fontsize{10}{12.5}\selectfont] at (0,24.7) {Bellingham Math Circle / Week 35 / W35-k-1-v2};
\node[anchor=north west,inner sep=0,align=left,text width=0.5cm,font=\sffamily\fontsize{10}{12.5}\selectfont] at (18,24.7) {1};
\node[anchor=north west,inner sep=0,align=left,text width=18.4cm,font=\sffamily\fontsize{14}{17.5}\selectfont] at (0,0.3) {Use the loose footprints and long strips. Each border repeats forever both ways. A match must fit every footprint. Only footprints count, not strip ends or guide lines. A slide must move the copy. Test with a tracing copy; leave the original in place.};
\node[anchor=north west,inner sep=0,align=left,text width=18.4cm,font=\sffamily\fontsize{17}{21.25}\selectfont] at (0,3.2) {\textbf{Problem 1:} Make one repeating border. Find slides that make its tracing copy match and slides that do not.};
\draw[gray!50,line width=.6pt] (.35,6.1) rectangle (18,12.1);
\draw[dashed,line width=.6pt] (.35,9.1) -- (18,9.1);
\draw[-{Stealth},line width=.9pt] (0.6,9.1) -- (0,9.1);
\draw[-{Stealth},line width=.9pt] (17.799999999999997,9.1) -- (18.4,9.1);
\draw[gray!45] (1.7,8.98) -- (1.7,9.219999999999999);
\draw[gray!45] (4.7,8.98) -- (4.7,9.219999999999999);
\draw[gray!45] (7.7,8.98) -- (7.7,9.219999999999999);
\draw[gray!45] (10.7,8.98) -- (10.7,9.219999999999999);
\draw[gray!45] (13.7,8.98) -- (13.7,9.219999999999999);
\draw[gray!45] (16.7,8.98) -- (16.7,9.219999999999999);\end{scope}\end{scope}\end{tikzpicture}
\newpage
\null\begin{tikzpicture}[remember picture,overlay]
\begin{scope}[shift={(current page.north west)},x=1cm,y=-1cm]
\begin{scope}[shift={(1.5,1.5)}]
\node[anchor=north west,inner sep=0,align=left,text width=18.5cm,font=\sffamily\fontsize{11}{13.75}\selectfont] at (0,-0.5) {Week 35 / Footprint borders / K--1};
\node[anchor=north west,inner sep=0,align=left,text width=17cm,font=\sffamily\fontsize{10}{12.5}\selectfont] at (0,24.7) {Bellingham Math Circle / Week 35 / W35-k-1-v2};
\node[anchor=north west,inner sep=0,align=left,text width=0.5cm,font=\sffamily\fontsize{10}{12.5}\selectfont] at (18,24.7) {2};
\node[anchor=north west,inner sep=0,align=left,text width=18.4cm,font=\sffamily\fontsize{17}{21.25}\selectfont] at (0,0.5) {\textbf{Problem 2:} Make two repeating borders whose shortest matching slides are different lengths.};
\draw[gray!50,line width=.6pt] (.35,4.1) rectangle (18,10.1);
\draw[dashed,line width=.6pt] (.35,7.1) -- (18,7.1);
\draw[-{Stealth},line width=.9pt] (0.6,7.1) -- (0,7.1);
\draw[-{Stealth},line width=.9pt] (17.799999999999997,7.1) -- (18.4,7.1);
\draw[gray!45] (1.7,6.9799999999999995) -- (1.7,7.22);
\draw[gray!45] (4.7,6.9799999999999995) -- (4.7,7.22);
\draw[gray!45] (7.7,6.9799999999999995) -- (7.7,7.22);
\draw[gray!45] (10.7,6.9799999999999995) -- (10.7,7.22);
\draw[gray!45] (13.7,6.9799999999999995) -- (13.7,7.22);
\draw[gray!45] (16.7,6.9799999999999995) -- (16.7,7.22);
\draw[gray!50,line width=.6pt] (.35,13.100000000000001) rectangle (18,19.1);
\draw[dashed,line width=.6pt] (.35,16.1) -- (18,16.1);
\draw[-{Stealth},line width=.9pt] (0.6,16.1) -- (0,16.1);
\draw[-{Stealth},line width=.9pt] (17.799999999999997,16.1) -- (18.4,16.1);
\draw[gray!45] (1.7,15.980000000000002) -- (1.7,16.220000000000002);
\draw[gray!45] (4.7,15.980000000000002) -- (4.7,16.220000000000002);
\draw[gray!45] (7.7,15.980000000000002) -- (7.7,16.220000000000002);
\draw[gray!45] (10.7,15.980000000000002) -- (10.7,16.220000000000002);
\draw[gray!45] (13.7,15.980000000000002) -- (13.7,16.220000000000002);
\draw[gray!45] (16.7,15.980000000000002) -- (16.7,16.220000000000002);\end{scope}\end{scope}\end{tikzpicture}
\newpage
\null\begin{tikzpicture}[remember picture,overlay]
\begin{scope}[shift={(current page.north west)},x=1cm,y=-1cm]
\begin{scope}[shift={(1.5,1.5)}]
\node[anchor=north west,inner sep=0,align=left,text width=18.5cm,font=\sffamily\fontsize{11}{13.75}\selectfont] at (0,-0.5) {Week 35 / Footprint borders / K--1};
\node[anchor=north west,inner sep=0,align=left,text width=17cm,font=\sffamily\fontsize{10}{12.5}\selectfont] at (0,24.7) {Bellingham Math Circle / Week 35 / W35-k-1-v2};
\node[anchor=north west,inner sep=0,align=left,text width=0.5cm,font=\sffamily\fontsize{10}{12.5}\selectfont] at (18,24.7) {3};
\begin{scope}[shift={(0,-2.5)}]
\draw[dashed,gray] (0.5,4.9) -- (4.7,4.9);
\draw[dashed,gray] (7.1,4.9) -- (11.299999999999999,4.9);
\draw[dashed,gray] (13.700000000000001,4.9) -- (17.900000000000002,4.9);
\begin{scope}[shift={(2.6,3.6)},rotate=0,xscale=0.68,yscale=0.68]
\draw[line width=.7pt,fill=gray!15] (-1.3,-1.1) -- (1.3,-1.1) -- (1.3,-.4) -- (.8,-.4) -- (.8,-.65) -- (-.5,-.65) -- (-.5,0) -- (.6,0) -- (.6,.45) -- (-.5,.45) -- (-.5,1.1) -- (-1.3,1.1) -- cycle;
\draw[-{Stealth},line width=.65pt] (-.25,-.86) -- (.55,-.86);\end{scope}
\begin{scope}[shift={(9.2,6.2)},rotate=0,xscale=0.68,yscale=-0.68]
\draw[line width=.7pt,fill=gray!15] (-1.3,-1.1) -- (1.3,-1.1) -- (1.3,-.4) -- (.8,-.4) -- (.8,-.65) -- (-.5,-.65) -- (-.5,0) -- (.6,0) -- (.6,.45) -- (-.5,.45) -- (-.5,1.1) -- (-1.3,1.1) -- cycle;
\draw[-{Stealth},line width=.65pt] (-.25,-.86) -- (.55,-.86);\end{scope}
\begin{scope}[shift={(16.8,6.2)},rotate=0,xscale=0.68,yscale=-0.68]
\draw[line width=.7pt,fill=gray!15] (-1.3,-1.1) -- (1.3,-1.1) -- (1.3,-.4) -- (.8,-.4) -- (.8,-.65) -- (-.5,-.65) -- (-.5,0) -- (.6,0) -- (.6,.45) -- (-.5,.45) -- (-.5,1.1) -- (-1.3,1.1) -- cycle;
\draw[-{Stealth},line width=.65pt] (-.25,-.86) -- (.55,-.86);\end{scope}
\node[anchor=north west,inner sep=0,align=left,text width=0.7cm,font=\sffamily\fontsize{12}{15.0}\selectfont] at (1.05,3.3000000000000003) {A};
\node[anchor=north west,inner sep=0,align=left,text width=0.7cm,font=\sffamily\fontsize{12}{15.0}\selectfont] at (7.6499999999999995,5.9) {A};
\node[anchor=north west,inner sep=0,align=left,text width=0.7cm,font=\sffamily\fontsize{12}{15.0}\selectfont] at (15.25,5.9) {A};
\draw[-{Stealth},line width=.9pt] (4.8,4.9) -- (6.8,4.9);
\draw[-{Stealth},line width=.9pt] (11.5,4.9) -- (13.4,4.9);
\node[anchor=north west,inner sep=0,align=left,text width=4.5cm,font=\sffamily\fontsize{12}{15.0}\selectfont] at (0.7,7.4) {start};
\node[anchor=north west,inner sep=0,align=left,text width=5cm,font=\sffamily\fontsize{12}{15.0}\selectfont] at (7.1,7.4) {flip over the line};
\node[anchor=north west,inner sep=0,align=left,text width=4.7cm,font=\sffamily\fontsize{12}{15.0}\selectfont] at (13.7,7.4) {then slide right};
\draw[-{Stealth},line width=.8pt] (14.7,7.1) -- (15.7,7.1);
\end{scope}
\node[anchor=north west,inner sep=0,align=left,text width=18.4cm,font=\sffamily\fontsize{17}{21.25}\selectfont] at (0,6.6) {\textbf{Problem 3:} Make a border that matches after a flip over the middle line and a slide along it. Can you make the flip alone fail?};
\draw[gray!50,line width=.6pt] (.35,9.3) rectangle (18,15.3);
\draw[dashed,line width=.6pt] (.35,12.3) -- (18,12.3);
\draw[-{Stealth},line width=.9pt] (0.6,12.3) -- (0,12.3);
\draw[-{Stealth},line width=.9pt] (17.799999999999997,12.3) -- (18.4,12.3);
\draw[gray!45] (1.7,12.180000000000001) -- (1.7,12.42);
\draw[gray!45] (4.7,12.180000000000001) -- (4.7,12.42);
\draw[gray!45] (7.7,12.180000000000001) -- (7.7,12.42);
\draw[gray!45] (10.7,12.180000000000001) -- (10.7,12.42);
\draw[gray!45] (13.7,12.180000000000001) -- (13.7,12.42);
\draw[gray!45] (16.7,12.180000000000001) -- (16.7,12.42);
\node[anchor=north west,inner sep=0,align=left,text width=18.4cm,font=\sffamily\fontsize{17}{21.25}\selectfont] at (0,16.3) {\textbf{Problem 4:} Make a different border with the same kind of match.};
\draw[gray!50,line width=.6pt] (.35,17.8) rectangle (18,23.8);
\draw[dashed,line width=.6pt] (.35,20.8) -- (18,20.8);
\draw[-{Stealth},line width=.9pt] (0.6,20.8) -- (0,20.8);
\draw[-{Stealth},line width=.9pt] (17.799999999999997,20.8) -- (18.4,20.8);
\draw[gray!45] (1.7,20.68) -- (1.7,20.92);
\draw[gray!45] (4.7,20.68) -- (4.7,20.92);
\draw[gray!45] (7.7,20.68) -- (7.7,20.92);
\draw[gray!45] (10.7,20.68) -- (10.7,20.92);
\draw[gray!45] (13.7,20.68) -- (13.7,20.92);
\draw[gray!45] (16.7,20.68) -- (16.7,20.92);\end{scope}\end{scope}\end{tikzpicture}
\newpage
\null\begin{tikzpicture}[remember picture,overlay]
\begin{scope}[shift={(current page.north west)},x=1cm,y=-1cm]
\begin{scope}[shift={(1.5,1.5)}]
\node[anchor=north west,inner sep=0,align=left,text width=18.5cm,font=\sffamily\fontsize{11}{13.75}\selectfont] at (0,-0.5) {Week 35 / Footprint borders / K--1};
\node[anchor=north west,inner sep=0,align=left,text width=17cm,font=\sffamily\fontsize{10}{12.5}\selectfont] at (0,24.7) {Bellingham Math Circle / Week 35 / W35-k-1-v2};
\node[anchor=north west,inner sep=0,align=left,text width=0.5cm,font=\sffamily\fontsize{10}{12.5}\selectfont] at (18,24.7) {4};
\node[anchor=north west,inner sep=0,align=left,text width=18.4cm,font=\sffamily\fontsize{17}{21.25}\selectfont] at (0,0.5) {\textbf{Problem 5:} Make a border that matches when you flip it over the middle line. Change it so that the flip fails.};
\draw[gray!50,line width=.6pt] (.35,4.1) rectangle (18,10.1);
\draw[dashed,line width=.6pt] (.35,7.1) -- (18,7.1);
\draw[-{Stealth},line width=.9pt] (0.6,7.1) -- (0,7.1);
\draw[-{Stealth},line width=.9pt] (17.799999999999997,7.1) -- (18.4,7.1);
\draw[gray!45] (1.7,6.9799999999999995) -- (1.7,7.22);
\draw[gray!45] (4.7,6.9799999999999995) -- (4.7,7.22);
\draw[gray!45] (7.7,6.9799999999999995) -- (7.7,7.22);
\draw[gray!45] (10.7,6.9799999999999995) -- (10.7,7.22);
\draw[gray!45] (13.7,6.9799999999999995) -- (13.7,7.22);
\draw[gray!45] (16.7,6.9799999999999995) -- (16.7,7.22);
\draw[gray!50,line width=.6pt] (.35,13.100000000000001) rectangle (18,19.1);
\draw[dashed,line width=.6pt] (.35,16.1) -- (18,16.1);
\draw[-{Stealth},line width=.9pt] (0.6,16.1) -- (0,16.1);
\draw[-{Stealth},line width=.9pt] (17.799999999999997,16.1) -- (18.4,16.1);
\draw[gray!45] (1.7,15.980000000000002) -- (1.7,16.220000000000002);
\draw[gray!45] (4.7,15.980000000000002) -- (4.7,16.220000000000002);
\draw[gray!45] (7.7,15.980000000000002) -- (7.7,16.220000000000002);
\draw[gray!45] (10.7,15.980000000000002) -- (10.7,16.220000000000002);
\draw[gray!45] (13.7,15.980000000000002) -- (13.7,16.220000000000002);
\draw[gray!45] (16.7,15.980000000000002) -- (16.7,16.220000000000002);\end{scope}\end{scope}\end{tikzpicture}
\newpage
\null\begin{tikzpicture}[remember picture,overlay]
\begin{scope}[shift={(current page.north west)},x=1cm,y=-1cm]
\begin{scope}[shift={(1.5,1.5)}]
\node[anchor=north west,inner sep=0,align=left,text width=18.5cm,font=\sffamily\fontsize{11}{13.75}\selectfont] at (0,-0.5) {Week 35 / Footprint borders / K--1};
\node[anchor=north west,inner sep=0,align=left,text width=17cm,font=\sffamily\fontsize{10}{12.5}\selectfont] at (0,24.7) {Bellingham Math Circle / Week 35 / W35-k-1-v2};
\node[anchor=north west,inner sep=0,align=left,text width=0.5cm,font=\sffamily\fontsize{10}{12.5}\selectfont] at (18,24.7) {5};
\node[anchor=north west,inner sep=0,align=left,text width=18.4cm,font=\sffamily\fontsize{14}{17.5}\selectfont] at (0,0.3) {A half-turn turns the copy halfway around a chosen point on the table. You may choose any point.};
\node[anchor=north west,inner sep=0,align=left,text width=18.4cm,font=\sffamily\fontsize{17}{21.25}\selectfont] at (0,2.1) {\textbf{Problem 6:} Make a border that matches after a half-turn. Make another that cannot match after a half-turn around any point.};
\draw[gray!50,line width=.6pt] (.35,5.300000000000001) rectangle (18,11.3);
\draw[dashed,line width=.6pt] (.35,8.3) -- (18,8.3);
\draw[-{Stealth},line width=.9pt] (0.6,8.3) -- (0,8.3);
\draw[-{Stealth},line width=.9pt] (17.799999999999997,8.3) -- (18.4,8.3);
\draw[gray!45] (1.7,8.180000000000001) -- (1.7,8.42);
\draw[gray!45] (4.7,8.180000000000001) -- (4.7,8.42);
\draw[gray!45] (7.7,8.180000000000001) -- (7.7,8.42);
\draw[gray!45] (10.7,8.180000000000001) -- (10.7,8.42);
\draw[gray!45] (13.7,8.180000000000001) -- (13.7,8.42);
\draw[gray!45] (16.7,8.180000000000001) -- (16.7,8.42);
\draw[gray!50,line width=.6pt] (.35,13.100000000000001) rectangle (18,19.1);
\draw[dashed,line width=.6pt] (.35,16.1) -- (18,16.1);
\draw[-{Stealth},line width=.9pt] (0.6,16.1) -- (0,16.1);
\draw[-{Stealth},line width=.9pt] (17.799999999999997,16.1) -- (18.4,16.1);
\draw[gray!45] (1.7,15.980000000000002) -- (1.7,16.220000000000002);
\draw[gray!45] (4.7,15.980000000000002) -- (4.7,16.220000000000002);
\draw[gray!45] (7.7,15.980000000000002) -- (7.7,16.220000000000002);
\draw[gray!45] (10.7,15.980000000000002) -- (10.7,16.220000000000002);
\draw[gray!45] (13.7,15.980000000000002) -- (13.7,16.220000000000002);
\draw[gray!45] (16.7,15.980000000000002) -- (16.7,16.220000000000002);
\node[anchor=north west,inner sep=0,align=left,text width=18.4cm,font=\sffamily\fontsize{17}{21.25}\selectfont] at (0,21.2) {\textbf{Problem 7:} Trade borders with your partner and find the moves that match.};\end{scope}\end{scope}\end{tikzpicture}
\end{document}
